\section{Fix precision first, then quantify}
\label{sec:fixedprecision}

\subsection{The quantifier-order problem}

Many packets describe a construction that must work at every requested
precision \(p\).  A tempting induction predicate is
\[
   \forall p\;\Inv(i,p,u).
\]
This is not quantifier free, so it cannot be fed to \code{PRA.qfInd}.  Even when
one can encode the universal quantifier by a bounded fold, the bound may depend
on the loop counter in an inconvenient way.  The correct proof order is:
\begin{enumerate}
\item take a single arbitrary precision \(p\) as a free parameter;
\item compute all finite queries needed at that precision;
\item perform quantifier-free induction with \(p\) free;
\item universally introduce \(p\) after the induction is complete.
\end{enumerate}

\subsection{Finite query tables}

Suppose a construction at output precision \(p\) makes finitely many oracle
queries to its input presentation.  Let
\[
   Q(p)=\langle q_0(p),\ldots,q_{m(p)-1}(p)\rangle
\]
be the primitive-recursive query table, and let
\(A(p)\) encode the corresponding finite answers.  The condition that all
answers are valid is the Boolean fold
\[
   \operatorname{TableOK}(x,p,A)
   =\operatorname{All}\bigl(m(p),
       j\mapsto \operatorname{AnswerOK}(x,q_j(p),A_j)\bigr).
\]
It is one quantifier-free equation.  The loop invariant may mention \(p,Q(p),A\)
freely without quantifying over other precisions.

\begin{theorem}[Fixed-precision induction]
Let \(\Phi(i,p,u)\) be a fixed bounded formula and let
\(b_\Phi(i,p,u)\) be its Booleanization.  If iPRA proves the base and successor
steps uniformly with free parameters \(p,u\), then iPRA proves
\[
      \forall p\,\forall u\,\forall i\;\Phi(i,p,u).
\]
The only induction used is quantifier-free induction on the equation
\(b_\Phi(i,p,u)=1\).
\end{theorem}

\begin{proof}
Work with \(p,u\) as free variables.  By \cref{thm:booleanization}, replace
\(\Phi\) by the equivalent equation \(b_\Phi=1\).  Apply quantifier-free
induction on \(i\), using the supplied base and successor derivations after
transport across the Booleanization equivalence.  This yields
\(\forall i\,b_\Phi(i,p,u)=1\) with \(p,u\) still free.  Convert back to
\(\Phi\), then use universal introduction first for \(u\) and then for \(p\).
The eigenvariable conditions hold because the context was shifted exactly as
required by the checked \code{Prf.allI} constructor.
\end{proof}

\subsection{Simultaneous finite windows}

Sometimes a theorem needs a whole window of precisions
\(p,p+1,\ldots,p+k\), for example to verify a tail estimate or a finite
coherence condition.  Do not introduce \(k+1\) independent quantified
invariants.  Extend the query table to contain the union of all queries needed
by the window and fold their validity into one Boolean.  The window length may
be a free parameter if its loop is bounded and compiled into the Boolean
state.

\subsection{Changing parameters in recursive algorithms}

An algorithm may refine a precision or alter a denominator cap at each step.
Let \(d(i,p,u)\) be the demand at stage \(i\), computed by primitive recursion.
The invariant is not ``all demands have been met'' as a quantified formula; it
is a Boolean statement about the finite state containing the answers already
collected for the current demand table.  A reverse-demand pass, described in
\cref{sec:dag}, computes a sufficiently strong initial table so that no new
unbounded oracle call is needed during the forward run.

\subsection{Examples}

\paragraph{Taylor approximation.}
Fix a requested error \(2^{-p}\).  Compute a derivative order \(N(p)\), a
radius margin, and all coefficient precisions before proving the finite Taylor
identity and remainder bound by induction on \(N(p)\).  Quantify over \(p\)
last.

\paragraph{Root clusters.}
Fix a disk radius and coefficient precision.  The finite root-count verifier
uses one grid and one finite boundary table.  It does not induct on a predicate
saying that the count is correct at every smaller radius.

\paragraph{Simplicial approximation.}
Fix the output subdivision level.  Compute the source subdivision and carrier
cloud required by the modulus.  Prove the chain-map and contiguity identities
for those finite complexes, then quantify over subdivision level.

\subsection{Why this resolves a recurring packet difficulty}

The difficulty reports often identify an induction that appears to need a
quantifier over all precisions.  In almost every such case, the quantifier is
an artifact of proving too strong a statement inside the loop.  The finite
computation only uses finitely many observations for a fixed requested
output.  Encoding exactly that finite-use table restores a quantifier-free
invariant and produces a theorem that is extensionally just as strong after
universal introduction.
